\FloatBarrier
\section{From Tunneling Amplitudes to Rate Equations}
\label{app:open-system-dynamics}

Section~\ref{sec:initialization-readout} uses reservoir entry and escape
rates to describe loading and spin-selective readout. This appendix
derives those rates from tunneling amplitudes and states the conditions
under which occupation probabilities obey a master equation.

\subsection{From tunneling amplitudes to a master equation}
\label{sec:amplitudes-to-master-equation}

An electron entering from the source and subsequently leaving
into the drain produces the charge cycle
\[
 \underbrace{\ket{N-1,0}}_{\ket A}
 \xrightarrow{\text{source}\to\text{dot}}
 \underbrace{\ket{N,0}}_{\ket B}
 \xrightarrow{\text{dot}\to\text{drain}}\ket A,
\]
where $0$ denotes a ground state. Quantum mechanics assigns
amplitudes to these tunneling processes. A rate equation instead
evolves the probabilities of occupying $A$ and $B$, treating entry
and exit as random events. We establish the connection between
these descriptions below.

\paragraph{Quantum amplitudes and the Green function.}
During an interval with fixed electrode voltages, write the
Hamiltonian of the dot and its reservoirs as
\begin{equation}
 H_{\mathrm{tot}}=H_{\mathrm{dot}}+H_{\mathrm R}+H_{\mathrm T},
 \qquad U(\tau)=e^{-iH_{\mathrm{tot}}\tau/\hbar},
 \label{eq:rates-total-propagator}
\end{equation}
where $H_{\mathrm T}$ transfers electrons between the dot and
the reservoirs. Let $\ket i,\ket j$ be dot states and
$\ket r,\ket s$ be many-electron states of the reservoirs;
$\ket{i,r}=\ket i\otimes\ket r$. The transition amplitude is
\begin{equation}
 \mathcal A_{js,ir}(\tau)=\bra{j,s}U(\tau)\ket{i,r}.
 \label{eq:rates-quantum-amplitude}
\end{equation}
It is a matrix element of the quantum propagator. With the
following normalization, the retarded Green operator for the
full Schr\"odinger equation is
\begin{equation}
 G^R(\tau)=-\frac{i}{\hbar}\theta(\tau)U(\tau),
 \qquad
 (i\hbar\partial_\tau-H_{\mathrm{tot}})G^R(\tau)
       =\delta(\tau)\mathbb I,
 \label{eq:rates-quantum-green-function}
\end{equation}
where $\theta$ is the step function. Thus the amplitude is related
to a matrix element of $G^R$, rather than being a probability.

\paragraph{Probabilities and the addition of channels.}
Prepare the dot in $\ket i$, initially uncorrelated with reservoirs
whose density matrix is
$\rho_{\mathrm R}=\sum_r w_r\ket r\bra r$, with
$w_r\geq0$ and $\sum_r w_r=1$. If only the final dot state is
measured, the conditional probability is
\begin{equation}
 P_{j\leftarrow i}(\tau)
   =\sum_{r,s}w_r\left|\mathcal A_{js,ir}(\tau)\right|^2.
 \label{eq:rates-conditional-probability}
\end{equation}
The sum over $r$ averages the initial reservoir mixture; the sum
over $s$ includes its unobserved, orthogonal final states.

For the two entry processes in
Appendix~\ref{app:dot-spectroscopy}, the final dot states
$\ket{N,0}$ and $\ket{N,1}$ are orthogonal. Suppressing the reservoir
labels for a moment, their contributions satisfy
\begin{equation}
 \left\|a_0\ket{N,0}+a_1\ket{N,1}\right\|^2
       =|a_0|^2+|a_1|^2.
 \label{eq:rates-orthogonal-channels}
\end{equation}
The same addition holds after summing over reservoir states.
Consequently, the probability of entry into either state is the
sum of the two entry probabilities. This follows from quantum
orthogonality, even if the final superposition is coherent.
By contrast, two coherent contributions $a^{(1)}$ and $a^{(2)}$
to the \emph{same full final state} give
$|a^{(1)}+a^{(2)}|^2$, including an interference term.

\paragraph{From a transition probability to a rate.}
To describe repeated tunneling by rates, assume weak coupling
to large reservoirs whose distributions remain approximately
fixed. Let $\tau_{\mathrm R}$ be the decay time of their correlation
functions. We consider intervals long enough for reservoir memory
to decay but short enough that a tunneling event is unlikely:
\begin{equation}
 \tau_{\mathrm R}\ll\Delta t\ll\Gamma_{\max}^{-1},
 \qquad
 P_{j\leftarrow i}(\Delta t)
       \simeq\Gamma_{j\leftarrow i}\Delta t\quad(j\ne i),
 \label{eq:rates-coarse-grained-definition}
\end{equation}
where $\Gamma_{\max}$ is the largest total escape rate from a
retained dot state. This defines a conditional rate: the dot is
already in $i$, and $\Gamma_{j\leftarrow i}$ gives its probability
per unit time of making the transition to $j$. Its unit is
$\mathrm{s}^{-1}$, not energy or velocity.

The linear dependence arises in weak-coupling perturbation theory
after summing over the reservoir continuum. To see the connection,
write $H_{\mathrm{dot}}\ket i=E_i\ket i$,
$H_{\mathrm R}\ket r=\epsilon_r\ket r$, and define
$T_{js,ir}=\bra{j,s}H_{\mathrm T}\ket{i,r}$ and
$\Delta E_{js,ir}=E_j+\epsilon_s-E_i-\epsilon_r$.
First-order time-dependent perturbation theory gives the
interaction-picture amplitude and its probability:
\begin{equation}
 \begin{aligned}
 \mathcal A^{(1)}_{js,ir}(\tau)
   &=-\frac{iT_{js,ir}}{\hbar}
       \int_0^\tau e^{i\Delta E_{js,ir}u/\hbar}\,du,\\
 |\mathcal A^{(1)}_{js,ir}(\tau)|^2
   &=\frac{4|T_{js,ir}|^2
       \sin^2(\Delta E_{js,ir}\tau/2\hbar)}{(\Delta E_{js,ir})^2}.
 \end{aligned}
 \label{eq:rates-perturbative-amplitude}
\end{equation}
Under a continuum energy integral, the long-time kernel obeys
\[
 \frac{4\sin^2(\Delta E\tau/2\hbar)}{\tau\Delta E^2}
     \longrightarrow\frac{2\pi}{\hbar}\delta(\Delta E).
\]
Substituting into
Eq.~\eqref{eq:rates-conditional-probability} gives Fermi's golden rule:
\begin{equation}
 \Gamma_{j\leftarrow i}
 =\frac{2\pi}{\hbar}\sum_{r,s}w_r|T_{js,ir}|^2
       \delta(E_j+\epsilon_s-E_i-\epsilon_r).
 \label{eq:rates-golden-rule}
\end{equation}
The delta function enforces conservation of the total energy.
The reservoir sum includes its available occupied and empty
states, giving the Fermi factors in
Eq.~\eqref{eq:measurement-tunnel-rates}.
Here the long-time form of the perturbative kernel is used within
the interval in Eq.~\eqref{eq:rates-coarse-grained-definition},
before appreciable depletion of the initial state. It is not
the exact $\tau\to0$ limit: at very short times the perturbative
probability above is proportional to $\tau^2$.

In the positive-bias regime where entry occurs from the source
and exit into the drain, set
$\Gamma_{\mathrm{in}}=\Gamma_{B\leftarrow A}$ and
$\Gamma_{\mathrm{out}}=\Gamma_{A\leftarrow B}$.
If two orthogonal final states can be loaded, probability addition
in Eq.~\eqref{eq:rates-orthogonal-channels} gives
\begin{equation}
 \Gamma_{\mathrm{in}}
   =\Gamma_{\mathrm S0}+\Gamma_{\mathrm S1}.
 \label{eq:rates-additive-entry}
\end{equation}
Each term is itself calculated from quantum amplitudes; adding
rates does not replace addition of amplitudes for coherent paths
to the same final state.

\paragraph{The master equation for occupation probabilities.}
A closed equation for occupations requires a further condition:
coherences between the retained dot states must not feed back
appreciably into their populations. For separated levels, rapidly
oscillating coherence terms can often be averaged out on the
same coarse time scale. Together with weak coupling and short
reservoir memory, this gives the population description used
here~\cite{Schaller2020Notes}. If coherences or reservoir memory
remain important, these occupation probabilities alone do not
provide a closed description.

Let $p(t)$ be the probability of $B$ and $1-p(t)$ that of $A$.
Over one coarse time step,
\begin{equation}
 \begin{aligned}
 p(t+\Delta t)\simeq{}&
 \underbrace{p(t)(1-\Gamma_{\mathrm{out}}\Delta t)}_
             {\text{already in }B\text{ and remains there}}\\
 &+\underbrace{[1-p(t)]\Gamma_{\mathrm{in}}\Delta t}_
             {\text{starts in }A\text{ and enters }B}.
 \end{aligned}
 \label{eq:rates-probability-step}
\end{equation}
Subtracting $p(t)$ and dividing by $\Delta t$ yields the differential
equation for the coarse-grained evolution,
\begin{equation}
 \dot p=\Gamma_{\mathrm{in}}(1-p)-\Gamma_{\mathrm{out}}p.
 \label{eq:rates-two-state-master-equation}
\end{equation}
This is the two-state \emph{Pauli master equation}, and reproduces
Eq.~\eqref{eq:measurement-zero-temperature-master-equation} with
$\Gamma_{\mathrm{in}}=\Gamma_{\mathrm S}$ and
$\Gamma_{\mathrm{out}}=\Gamma_{\mathrm D}$.
For more dot states, with occupation probabilities $p_i$,
the same gain-minus-loss argument gives
\begin{equation}
 \dot p_j=\sum_{i\ne j}\Gamma_{j\leftarrow i}p_i
           -p_j\sum_{i\ne j}\Gamma_{i\leftarrow j}.
 \label{eq:rates-general-master-equation}
\end{equation}
Orthogonal-channel addition alone does not justify merging their
final-state occupations. The two-state treatment of excited-state
transport in Appendix~\ref{app:dot-spectroscopy} also
uses fast relaxation to the $N$-electron ground state; otherwise
the excited states need their own $p_i$.

\paragraph{Waiting times and current.}
For a constant total entry rate, the probability $q(t)$ that no
electron has yet entered, starting in $A$, obeys
$q(t+\Delta t)\simeq q(t)(1-\Gamma_{\mathrm{in}}\Delta t)$.
Thus
\begin{equation}
 q(t)=e^{-\Gamma_{\mathrm{in}}t},\qquad
 \langle\tau_{\mathrm{in}}\rangle
   =\int_0^\infty q(t)\,dt=\Gamma_{\mathrm{in}}^{-1}.
 \label{eq:rates-waiting-time}
\end{equation}
The corresponding mean waiting time for exit, starting in $B$,
is $\Gamma_{\mathrm{out}}^{-1}$. These are conditional waiting
times. In a running transport cycle, the ensemble-averaged numbers
of entries and exits per unit time are instead
$\Gamma_{\mathrm{in}}(1-p)$ and $\Gamma_{\mathrm{out}}p$.
In the stationary state these are equal, and the current convention
of Appendix~\ref{app:charge-detection} gives
\begin{equation}
 p_{\mathrm{st}}=\frac{\Gamma_{\mathrm{in}}}
                         {\Gamma_{\mathrm{in}}+\Gamma_{\mathrm{out}}},
 \qquad
 I=e\Gamma_{\mathrm{out}}p_{\mathrm{st}}
   =e\frac{\Gamma_{\mathrm{in}}\Gamma_{\mathrm{out}}}
                 {\Gamma_{\mathrm{in}}+\Gamma_{\mathrm{out}}}.
 \label{eq:rates-stationary-current}
\end{equation}
The rates are therefore quantum-mechanically determined inputs
to a classical stochastic approximation for the occupations;
neither a rate nor an occupation probability is a quantum
propagation amplitude.
